%=======================================================================
%  Provably Secure Four-Phase PUF-Based Authentication and Key
%  Establishment for UAV Swarms
%
%  Single-file, self-contained manuscript.
%  Compile:  pdflatex uav_puf_protocol.tex   (run twice for references)
%=======================================================================
\documentclass[11pt,a4paper]{article}

% --- Page geometry -----------------------------------------------------
\usepackage[a4paper,margin=1in]{geometry}

% --- Fonts & encoding --------------------------------------------------
\usepackage[T1]{fontenc}
\usepackage{lmodern}

% --- Mathematics -------------------------------------------------------
\usepackage{amsmath,amssymb,amsfonts,amsthm}
\usepackage{mathtools}
\usepackage{bm}
\allowdisplaybreaks
\setlength{\jot}{7pt}          % generous spacing between aligned rows

% --- Tables / lists ----------------------------------------------------
\usepackage{booktabs}
\usepackage{array}
\usepackage{tabularx}
\usepackage{enumitem}
\setlist{itemsep=0.5em,topsep=0.4em}

% --- Colour (define before hyperref) -----------------------------------
\usepackage{xcolor}
\definecolor{navy}{RGB}{22,45,95}
\definecolor{steel}{RGB}{60,80,120}
\definecolor{boxbg}{RGB}{245,247,251}
\definecolor{boxln}{RGB}{200,210,228}

% --- Algorithms --------------------------------------------------------
\usepackage{algorithm}
\usepackage{algorithmic}
\usepackage{setspace}
\renewcommand{\algorithmicrequire}{\textbf{Input:}}
\renewcommand{\algorithmicensure}{\textbf{Output:}}

% --- Framed boxes ------------------------------------------------------
\usepackage{mdframed}

% --- Section styling ---------------------------------------------------
\usepackage{titlesec}
\titleformat{\section}{\normalfont\Large\bfseries\color{navy}}{\thesection.}{0.6em}{}
\titleformat{\subsection}{\normalfont\large\bfseries\color{steel}}{\thesubsection}{0.6em}{}
\titlespacing*{\section}{0pt}{1.6em plus .3em}{0.75em}
\titlespacing*{\subsection}{0pt}{1.15em plus .2em}{0.5em}

% --- Microtypography ---------------------------------------------------
\usepackage{microtype}

% --- Hyperlinks (before cleveref) --------------------------------------
\usepackage[colorlinks=true,linkcolor=navy,citecolor=navy,urlcolor=navy]{hyperref}
\usepackage[capitalize]{cleveref}

% --- Paragraph look ----------------------------------------------------
\setlength{\parindent}{0pt}
\setlength{\parskip}{0.55em plus 0.1em}

% --- Theorem environments (spaced, coloured heads) ---------------------
\newtheoremstyle{jrnlplain}{1.2em}{1.2em}{\itshape}{}{\bfseries\color{navy}}{.}{.6em}{}
\newtheoremstyle{jrnldef}{1.2em}{1.2em}{}{}{\bfseries\color{navy}}{.}{.6em}{}
\theoremstyle{jrnlplain}
\newtheorem{theorem}{Theorem}
\newtheorem{lemma}[theorem]{Lemma}
\newtheorem{corollary}[theorem]{Corollary}
\theoremstyle{jrnldef}
\newtheorem{definition}{Definition}
\newtheorem{assumption}{Assumption}
\newtheorem{remark}{Remark}

% --- Notation macros ---------------------------------------------------
\newcommand{\GS}{\mathrm{GS}}
\newcommand{\UAV}{\mathrm{UAV}}
\newcommand{\PUF}{\mathsf{PUF}}
\newcommand{\TID}{\mathrm{TID}}
\newcommand{\Cred}{\mathrm{Cred}}
\newcommand{\mk}{\mathrm{mk}}
\newcommand{\SK}{\mathrm{SK}}
\newcommand{\MAC}{\mathsf{MAC}}
\newcommand{\aead}{\mathsf{AE}}
\newcommand{\Enc}{\mathsf{Enc}}
\newcommand{\Dec}{\mathsf{Dec}}
\newcommand{\KDF}{\mathsf{KDF}}
\newcommand{\Ext}{\mathsf{Ext}}
\newcommand{\Codec}{\mathsf{C}}
\newcommand{\Gen}{\mathsf{Gen}}
\newcommand{\Rep}{\mathsf{Rep}}
\newcommand{\FE}{\mathsf{FE}}
\newcommand{\adv}{\mathcal{A}}
\newcommand{\secpar}{\lambda}
\newcommand{\bits}{\{0,1\}}
\newcommand{\getsr}{\xleftarrow{\$}}
\newcommand{\negl}{\mathsf{negl}}
\newcommand{\Adv}{\mathsf{Adv}}
\newcommand{\entropy}{\widetilde{H}_\infty}
\newcommand{\lbl}[1]{\text{``}\texttt{#1}\text{''}}
\newcommand{\catt}{\,\|\,}
\newcommand{\xorop}{\oplus}
\newcommand{\transcript}{\mathcal{T}}
\newcommand{\cmnt}[1]{\hfill\ensuremath{\triangleright}~\textit{#1}}

\title{\vspace{-1.2em}\bfseries\color{navy}
Provably Secure Four-Phase PUF-Based Authentication\\[2pt]
and Key Establishment for UAV Swarms}
\author{Md~Ziyad~Hussain \qquad Rakesh~Matam\\[2pt]
\normalsize Department of Computer Science and Engineering\\[1pt]
\normalsize Indian Institute of Information Technology Guwahati, Assam, India}
\date{}

\begin{document}
\maketitle

%-----------------------------------------------------------------------
\begin{abstract}
\noindent
Unmanned Aerial Vehicle (UAV) swarms require authentication that is fast,
lightweight, and resilient to the physical capture of individual drones.
We present a four-phase authentication and key-establishment protocol whose
root of trust is a Physical Unclonable Function (PUF) embedded in each UAV.
The protocol (i)~derives a \emph{stable} per-device master key from the PUF
through a fuzzy extractor with an exact min-entropy budget, so that no secret
key is ever stored on the drone; (ii)~authenticates the UAV and the Ground
Station (GS) to one another with keyed message-authentication codes, keeping
the raw PUF response entirely on-device; (iii)~establishes GS-free,
\emph{per-pair} peer credentials delivered under authenticated encryption;
and (iv)~binds a mandatory, authenticated ephemeral Diffie--Hellman exchange
into every session key to guarantee perfect forward secrecy. We prove six
security theorems---mutual authentication in both directions, peer mutual
authentication, replay/MITM resistance, physical-capture resistance,
session-key secrecy, and perfect forward secrecy---each by a game-based
reduction to standard assumptions, and we corroborate the authentication and
secrecy claims with a symbolic model in the Tamarin prover. The design uses
only a $160$-bit hash (instantiable as SHA3-160 or SPONGENT-160), a keyed
MAC, a lightweight AEAD, a key-derivation function, and a single elliptic
curve, making it well suited to resource-constrained aerial platforms.
\end{abstract}

\medskip

%=======================================================================
\section{Introduction}
%=======================================================================
Multi-UAV swarms coordinate continuously with a Ground Station over a
command-and-control link and with one another over peer links to maintain
formation, fuse sensor data, and relay traffic. Every such interaction
presupposes \emph{mutual authentication}: a UAV must verify the identity of
its GS and of each peer before accepting commands, sharing telemetry, or
forwarding messages. Public-key handshakes add latency and, more importantly,
depend on private keys held in non-volatile memory (NVM) that a captured drone
can surrender; pre-shared keys collapse the entire swarm once any single
device is compromised.

Physical Unclonable Functions offer a different foundation. A PUF maps
challenges to responses using nanometre-scale manufacturing variations unique
to one integrated circuit~\cite{suh2007}; it is unclonable, unpredictable, and
tamper-evident, and it stores no key---the ``key'' is the silicon itself.
Because PUF responses are noisy, a fuzzy extractor~\cite{dodis2008} reconciles
a field response to its enrolment value and distils a reproducible key.

\paragraph{Contributions.}
\begin{enumerate}[leftmargin=1.6em,label=\textbf{C\arabic*.}]
\item A four-phase protocol---secure enrolment, UAV--GS mutual
authenticated key exchange, GS-free UAV--UAV peer authentication, and
forward-secret session keying---built entirely from symmetric primitives plus
one elliptic curve.
\item A trust model in which each UAV holds \emph{no} stored secret: a stable
master key is regenerated on demand from the PUF and public helper data, with
an explicit min-entropy budget (\cref{lem:fe}) governing parameter choice.
\item Keyed authenticators throughout the handshake, so that the raw PUF
response never appears on the wire and cannot be harvested
(\cref{thm:p2,cor:nooracle}).
\item Per-pair peer credentials distributed under authenticated encryption,
confining the effect of a captured node to its own incident links
(\cref{thm:capture}), together with a mandatory authenticated ephemeral
exchange that delivers perfect forward secrecy (\cref{thm:pfs}).
\item A complete security analysis: six theorems with game-based reductions
to standard assumptions, mirrored by a mechanised Tamarin model
(\cref{sec:proofs,sec:symbolic}).
\end{enumerate}

\Cref{sec:prelim} fixes primitives and assumptions; \cref{sec:model} states
the system and adversary model; \cref{sec:design} gives the design
principles; \cref{sec:phases} specifies the four phases;
\cref{sec:proofs,sec:symbolic} present the computational and symbolic
analyses; \cref{sec:conclusion} concludes.

%=======================================================================
\section{Preliminaries and Cryptographic Primitives}
\label{sec:prelim}
%=======================================================================
Let $\secpar$ denote the security parameter. A quantity is \emph{negligible},
$\negl(\secpar)$, if it decays faster than any inverse polynomial. We write
$x\getsr S$ for uniform sampling, $\catt$ for concatenation (all
concatenations are length-prefixed and uniquely parseable), and $\xorop$ for
bitwise \textsc{xor}.

\subsection{Hash function}
$h:\bits^\ast\to\bits^{160}$ is a cryptographic hash instantiated in either of
two interoperable modes---\textbf{SHA3-256 truncated to $160$ bits}%
~\cite{fips202} or \textbf{SPONGENT-160}~\cite{bogdanov2011}---so that a
software-optimised and a hardware-lightweight primitive share one interface.
In the computational analysis $h$ is modelled as a random oracle. Keyed
functionality is provided by the dedicated constructions below rather than by
bare $h$.

\medskip
\subsection{Keyed primitives}

\begin{definition}[Message authentication code]
$\MAC:\bits^{\secpar}\times\bits^\ast\to\bits^{\secpar}$ is
\emph{existentially unforgeable under chosen-message attack} (EUF-CMA) if, for
$K\getsr\bits^{\secpar}$ and every probabilistic polynomial-time (PPT)
adversary $\adv$,
\begin{equation}
\Adv^{\mathrm{euf}}_{\MAC}(\adv)
   \;=\; \Pr\!\big[\,\adv\ \text{outputs a fresh valid}\ (m,\MAC_K(m))\,\big]
   \;\le\; \negl(\secpar).
\end{equation}
We instantiate $\MAC$ as $\mathrm{HMAC}$ over $h$ (equivalently, a keyed
duplex/sponge MAC).
\end{definition}

\begin{definition}[Authenticated encryption]
A nonce-based AEAD scheme $\aead=(\Enc,\Dec)$ computes
\begin{align}
c &= \Enc_K(\eta;\ \mathit{ad};\ m),
   &
m/\!\bot &= \Dec_K(\eta;\ \mathit{ad};\ c),
\end{align}
and provides ciphertext indistinguishability (IND-CCA) and ciphertext
integrity (INT-CTXT), with
$\Adv^{\mathrm{ind}}_{\aead},\Adv^{\mathrm{int}}_{\aead}\le\negl(\secpar)$ for
distinct nonces $\eta$. We use a lightweight nonce-based AEAD of the
Ascon family~\cite{ascon}, matching the resource profile of the platform.
\end{definition}

\begin{definition}[Key-derivation function]
$\KDF$ is HKDF (extract-then-expand)~\cite{krawczyk2010}. Each invocation
carries a unique \emph{domain-separation label}; $\KDF(k;\ell)$ denotes
expansion of key material $k$ under a label
\begin{equation}
\ell\in\{\lbl{p2-auth},\ \lbl{p2-sess},\ \lbl{pair-cred},\
         \lbl{peer-auth},\ \lbl{peer-sess}\}.
\end{equation}
We model $\KDF(\cdot;\ell)$ as an independent pseudorandom function (PRF) for
each label, so that distinct labels yield independent keys.
\end{definition}

\medskip
\subsection{Ephemeral key exchange}
Let $\mathbb{G}$ be the group underlying \textbf{X25519}~\cite{bernstein2006}
with generator $g$. Forward secrecy is reduced to the \emph{gap computational
Diffie--Hellman} (gap-CDH) assumption, equivalently the PRF-ODH assumption for
$\KDF$ over $\mathbb{G}$.

\medskip
\subsection{PUF and fuzzy extractor}
Each device $i$ embeds a strong PUF $\PUF_i:\bits^{c}\to\bits^{n}$. Two
evaluations of the same challenge differ in at most $t$ bits with
overwhelming probability.

\begin{assumption}[PUF unpredictability]\label{ass:puf}
For a challenge $C$ not previously queried, and given any polynomial set of
challenge--response pairs of $\PUF_i$, the response $\PUF_i(C)$ is
computationally indistinguishable from uniform to a PPT adversary that does
not physically possess device~$i$.
\end{assumption}

\begin{assumption}[Tamper evidence]\label{ass:tamper}
Invasive probing of device $i$ destroys $\PUF_i$; its challenge--response
behaviour can neither be extracted from captured silicon nor transferred to
another device.
\end{assumption}

We employ a fuzzy extractor $\FE=(\Gen,\Rep)$ formed from a code-offset secure
sketch over a binary linear code $\Codec$ with parameters $(n,k,t)$
(instantiated as $\mathrm{BCH}(255,131,18)$, hence $n-k=124$ parity bits and
correction radius $t=18$), composed with a strong extractor realised by
$\KDF$:
\begin{align}
\Gen(R):\quad
   & r\getsr\bits^{k},\qquad c=\Codec.\mathsf{Encode}(r),\qquad
     s = R\xorop c,\qquad \mathrm{seed}\getsr\bits^{\secpar},\\[1ex]
   & \mk=\KDF\big(\Ext_{\mathrm{seed}}(R);\ \lbl{root}\big),\qquad
     P=(s,\mathrm{seed}),\qquad \textbf{return } (P,\mk); \label{eq:gen}\\[2ex]
\Rep(R',P):\quad
   & c' = R'\xorop s,\qquad \hat r=\Codec.\mathsf{Decode}(c'),\qquad
     \hat c=\Codec.\mathsf{Encode}(\hat r),\qquad \hat R=\hat c\xorop s,\\[1ex]
   & \textbf{return } \mk=\KDF\big(\Ext_{\mathrm{seed}}(\hat R);\ \lbl{root}\big).
     \label{eq:rep}
\end{align}

\begin{lemma}[Reproduction correctness]\label{lem:fecorrect}
If the Hamming distance satisfies $d_H(R,R')\le t$, then $\Rep(R',P)$ recovers
$\hat R=R$ and hence the identical master key $\mk$.
\end{lemma}
\begin{proof}
$c'=R'\xorop s = c\xorop(R\xorop R')$ differs from the codeword $c$ in
$d_H(R,R')\le t$ positions, so bounded-distance decoding returns $\hat r=r$;
thus $\hat c=c$ and $\hat R=\hat c\xorop s = R$. Both parties therefore apply
$\Ext$ and $\KDF$ to the same value.
\end{proof}

\begin{lemma}[Helper-data privacy]\label{lem:fe}
Let $\entropy(R)=m$ be the conditional min-entropy of the enrolment response.
The code-offset sketch $s$ discloses at most $n-k$ bits:
\begin{equation}
\entropy\!\big(R \mid P\big)\ \ge\ m-(n-k).
\label{eq:leak}
\end{equation}
Consequently, if $\Ext$ is a $\big(m-(n-k),\varepsilon\big)$-strong extractor
with $\varepsilon=2^{-\Omega(\secpar)}$ and output length
$\ell\le m-(n-k)-2\log(1/\varepsilon)$, then $\mk$ is within statistical
distance $\varepsilon$ of uniform given $P$.
\end{lemma}
\begin{proof}
Given the internal randomness $r$, $s=R\xorop c$ is a deterministic function
of $R$ taking values in a coset of $\Codec$; revealing the coset of an
$[n,k]$ code fixes at most $n-k$ syndrome bits, so by the chain rule for
average min-entropy $\entropy(R\mid s)\ge\entropy(R)-(n-k)$, and
$\mathrm{seed}$ is independent of $R$. The strong-extractor property applied to
the residual source yields the claimed distance.
\end{proof}

\begin{remark}[Parameter selection]
Because a code-offset sketch is a bounded-leakage primitive
(\cref{eq:leak}), the enrolment aggregates $L$ independent response blocks
$R=(R^{(1)},\dots,R^{(L)})$ chosen so that the residual entropy meets the key
budget:
\begin{equation}
\sum_{b=1}^{L}\big(m_b-(n-k)\big)\ \ge\ \ell + 2\log(1/\varepsilon),
\qquad\text{e.g. } \ell=2\secpar,
\label{eq:design}
\end{equation}
where $m_b$ is the measured per-block min-entropy. This inequality determines
$L$ for a given PUF and target key length.
\end{remark}

\medskip
\begin{table}[t]
\centering
\caption{Notation.}
\label{tab:notation}
\renewcommand{\arraystretch}{1.25}
\begin{tabularx}{\linewidth}{@{}l X@{}}
\toprule
\textbf{Symbol} & \textbf{Meaning}\\
\midrule
$\GS$, $\UAV_i$ & ground station; the $i$-th UAV\\
$\TID_i$ & rotatable temporary identity of $\UAV_i$\\
$\PUF_i,\ R,R'$ & PUF of device $i$; enrolment / field response\\
$\FE=(\Gen,\Rep)$ & fuzzy extractor, \cref{eq:gen,eq:rep}\\
$P_i=(s_i,\mathrm{seed}_i)$ & public helper data\\
$\mk_i$ & stable per-device master key (never stored on the UAV)\\
$\mk_{\GS}$ & GS-only master key for pairwise credentials\\
$\Cred_{ij}$ & per-pair peer credential\\
$K^{i}_{\mathrm{auth}},\,K^{ij}_{\mathrm{auth}}$ & derived MAC keys\\
$\SK_i,\,\SK_{ij}$ & forward-secret session keys (Phase~2 / Phase~3)\\
$N_x,\ T_x$ & $128$-bit nonce; timestamp\\
$a,b;\ g^{a},g^{ab}$ & ephemeral scalars; public shares; shared secret\\
$\transcript$ & session transcript (all preceding message fields)\\
$\MAC,\aead,\KDF$ & keyed MAC; AEAD; key-derivation function\\
\bottomrule
\end{tabularx}
\end{table}

%=======================================================================
\section{System and Adversary Model}
\label{sec:model}
%=======================================================================
The system comprises one trusted $\GS$ with tamper-resistant storage and $N$
UAVs, each equipped with a PUF, a lightweight processor, and non-volatile
memory. We consider a strong adversary that combines Dolev--Yao network
control~\cite{dolevyao1983} with physical capture and an active-challenger
capability:

\begin{itemize}[leftmargin=1.5em]
\item \textbf{Network ($\adv_{\mathrm{net}}$).} Reads, injects, modifies,
reorders, replays, delays, and drops any message on any link.
\item \textbf{Active challenger ($\adv_{\mathrm{chal}}$).} May act as $\GS$
toward an honest UAV and submit adversarially chosen challenges and nonces.
\item \textbf{Capture ($\adv_{\mathrm{cap}}$).} May seize a device $k$ and
read its NVM contents $\{\TID_k,\,P_k,\,\{\Cred_{kj}\}_j\}$ together with any
live session keys. By \cref{ass:tamper} it obtains neither $\PUF_k$ nor the
transient master key $\mk_k$ (which is never stored).
\item \textbf{Post-session reveal.} For forward-secrecy analysis, the
adversary may learn long-term secrets ($\mk_i$, $\Cred_{ij}$) through a
\textsf{Reveal} query issued only \emph{after} the target session closes.
\end{itemize}

The adversary is PPT and, except with negligible probability, cannot invert
$h$, forge $\MAC$ or $\aead$, violate PUF unpredictability
(\cref{ass:puf}), or solve gap-CDH. Physical side channels are out of scope;
constant-time primitives are assumed.

%=======================================================================
\section{Design Principles}
\label{sec:design}
%=======================================================================
The protocol rests on five principles that shape every phase.

\begin{enumerate}[leftmargin=1.6em,label=\textbf{P\arabic*.}]
\item \textbf{Stable PUF-derived key, no stored secret.} Enrolment fixes a
stable master key $\mk_i$ with the fuzzy extractor. The UAV regenerates
$\mk_i$ on demand from its PUF and \emph{public} helper data; no secret resides
in NVM. Every authenticator is a MAC under a key derived from $\mk_i$, which
only the genuine device (through its PUF) and the $\GS$ (through its database)
can compute.

\item \textbf{The response stays on-device.} The fuzzy extractor executes
entirely on the UAV; no PUF response, and no function of a response under an
externally chosen mask, is ever transmitted.

\item \textbf{Verify before acting.} A UAV validates the $\GS$'s
authenticator before performing any PUF-dependent computation or emitting any
message, so a spurious challenger elicits no observable response.

\item \textbf{Per-pair credentials under AEAD.} Peer trust uses per-pair
credentials $\Cred_{ij}$ delivered confidentially and authentically under a
PUF-derived session key, with domain-separated labels and per-message-type
tags throughout.

\item \textbf{Mandatory authenticated forward secrecy.} Every session key
mixes in a fresh ephemeral Diffie--Hellman secret whose public shares are
bound into the MAC transcript; there is no non-ephemeral derivation path.
\end{enumerate}

%=======================================================================
\section{The Four-Phase Protocol}
\label{sec:phases}
%=======================================================================

\subsection{Phase 1: Secure Enrolment}
\label{sec:p1}
Enrolment is performed once per device over a secure out-of-band channel
(\cref{alg:p1}). The master key $\mk_i$ is the sole long-term secret and
resides only in the GS's tamper-resistant database; the UAV reconstructs it
transiently whenever needed.

\begin{algorithm}[t]
\caption{Secure Enrolment of $\UAV_i$}
\label{alg:p1}
\begin{spacing}{1.3}
\begin{algorithmic}[1]
\REQUIRE device $\UAV_i$ with embedded $\PUF_i$; target key length $\ell$
\ENSURE GS record $\langle\TID_i,\bm{C}_i,P_i,\mk_i\rangle$; UAV state $\langle\TID_i,\bm{C}_i,P_i\rangle$
\STATE $\GS$ selects a challenge block $\bm{C}_i=(C_i^{(1)},\dots,C_i^{(L)})$ and evaluates $R\gets\PUF_i(\bm{C}_i)$ in situ
\STATE $(P_i,\mk_i)\gets\FE.\Gen(R)$ \cmnt{$L$ chosen to satisfy the entropy budget~\eqref{eq:design}}
\STATE $\GS$ stores $\langle\TID_i,\bm{C}_i,P_i,\mk_i\rangle$ in tamper-resistant memory
\STATE $\UAV_i$ stores $\langle\TID_i,\bm{C}_i,P_i\rangle$ only \cmnt{$P_i$ is public; no secret resides in NVM}
\end{algorithmic}
\end{spacing}
\end{algorithm}

\subsection{Phase 2: UAV--GS Mutual Authenticated Key Exchange}
\label{sec:p2}
Phase~2 (\cref{alg:p2}) is a three-message mutual authenticated key exchange
under $K^{i}_{\mathrm{auth}}=\KDF(\mk_i;\lbl{p2-auth})$, with ephemeral shares
$g^{a}$ (UAV) and $g^{b}$ (GS) supplying forward secrecy.

\begin{algorithm}[t]
\caption{Phase 2 --- Mutual authentication and forward-secret keying}
\label{alg:p2}
\begin{spacing}{1.3}
\begin{algorithmic}[1]
\STATE $\UAV_i\!\to\!\GS$ \ $(M_1)$: sample $N_1,a$; \ $\sigma_1=\MAC_{K^{i}_{\mathrm{auth}}}(\lbl{M1}\catt\TID_i\catt N_1\catt T_1\catt g^{a})$; send $\langle\TID_i,N_1,T_1,g^{a},\sigma_1\rangle$
\STATE $\GS$: reject if $|T_{\mathrm{now}}-T_1|>\Delta T$ or $N_1$ already seen; fetch $\mk_i$; verify $\sigma_1$
\STATE $\GS\!\to\!\UAV_i$ \ $(M_2)$: sample $N_2,b$; \ $\sigma_2=\MAC_{K^{i}_{\mathrm{auth}}}(\lbl{M2}\catt\TID_i\catt N_1\catt N_2\catt T_2\catt g^{a}\catt g^{b})$; send $\langle N_2,T_2,g^{b},\sigma_2\rangle$
\STATE $\UAV_i$: recompute $\mk_i\gets\FE.\Rep(\PUF_i(\bm{C}_i),P_i)$; verify $\sigma_2$ \emph{before proceeding} \cmnt{a failed check aborts silently}
\STATE $\UAV_i\!\to\!\GS$ \ $(M_3)$: $\sigma_3=\MAC_{K^{i}_{\mathrm{auth}}}(\lbl{M3}\catt\transcript)$; send $\langle\sigma_3\rangle$
\STATE both set $\SK_i=\KDF(\mk_i\catt g^{ab}\catt\transcript;\ \lbl{p2-sess})$; $\GS$ rotates $\TID_i$
\end{algorithmic}
\end{spacing}
\end{algorithm}

\begin{remark}
No PUF response or masked response appears in any message; $R$ is consumed
locally by $\FE.\Rep$ in~\eqref{eq:rep}, and the $\GS$ already holds $\mk_i$.
Because $\mk_i$ is reusable, a single enrolment supports arbitrarily many
authentications; the challenge block $\bm{C}_i$ serves only periodic key
refresh, not per-session consumption.
\end{remark}

\subsection{Phase 3: UAV--UAV Peer Authentication}
\label{sec:p3}
After Phase~2, the $\GS$ mints per-pair credentials and distributes them under
authenticated encryption; peers then authenticate without further GS
involvement (\cref{alg:p3}).

\begin{algorithm}[t]
\caption{Phase 3 --- Per-pair credential issuance and peer handshake}
\label{alg:p3}
\begin{spacing}{1.3}
\begin{algorithmic}[1]
\STATE \textbf{Issuance.} for peers $i<j$: $\Cred_{ij}=\KDF(\mk_{\GS};\lbl{pair-cred}\catt i\catt j)$; \ $\GS\!\to\!\UAV_i$: $\aead.\Enc_{\SK_i}(\eta;\ i\catt j;\ \Cred_{ij})$, and symmetrically to $\UAV_j$
\STATE $\UAV_i$: $\Cred_{ij}\gets\aead.\Dec_{\SK_i}(\cdots)$; abort on $\bot$ \cmnt{confidential and authenticated}
\STATE set $K^{ij}_{\mathrm{auth}}=\KDF(\Cred_{ij};\lbl{peer-auth})$
\STATE $\UAV_i\!\to\!\UAV_j$ \ $(P_1)$: sample $n_i,a'$; \ $\tau_1=\MAC_{K^{ij}_{\mathrm{auth}}}(\lbl{P1}\catt i\catt j\catt n_i\catt T\catt g^{a'})$
\STATE $\UAV_j\!\to\!\UAV_i$ \ $(P_2)$: verify $\tau_1$; sample $n_j,b'$; \ $\tau_2=\MAC_{K^{ij}_{\mathrm{auth}}}(\lbl{P2}\catt i\catt j\catt n_i\catt n_j\catt g^{a'}\catt g^{b'})$
\STATE $\UAV_i\!\to\!\UAV_j$ \ $(P_3)$: verify $\tau_2$; \ $\tau_3=\MAC_{K^{ij}_{\mathrm{auth}}}(\lbl{P3}\catt\transcript)$
\STATE both set $\SK_{ij}=\KDF(\Cred_{ij}\catt g^{a'b'}\catt\transcript;\ \lbl{peer-sess})$
\end{algorithmic}
\end{spacing}
\end{algorithm}

\begin{remark}
Each message type carries a distinct tag ($\lbl{P1},\lbl{P2},\lbl{P3}$) and a
freshness element ($n_i,n_j,T$), and the session key uses label
$\lbl{peer-sess}$ independently of the authenticator key
$K^{ij}_{\mathrm{auth}}$ (label $\lbl{peer-auth}$), so the key and the
verifier are cryptographically separated.
\end{remark}

\subsection{Phase 4: Session Keys and Forward Secrecy}
\label{sec:p4}
Phase~4 is the forward-secret keying embedded in \cref{alg:p2,alg:p3}. Both
session keys share the unified form

\begin{mdframed}[backgroundcolor=boxbg,linecolor=boxln,linewidth=0.8pt,
  roundcorner=4pt,innertopmargin=8pt,innerbottommargin=8pt]
\begin{equation}
\SK \;=\; \KDF\Big(\,
   \underbrace{k_{\mathrm{lt}}}_{\mk_i\ \text{or}\ \Cred_{ij}}
   \ \catt\ g^{xy}\ \catt\ \transcript\ ;\ \ell \,\Big),
\label{eq:sk}
\end{equation}
\end{mdframed}

\noindent
where the ephemeral public shares $g^{x},g^{y}$ are authenticated inside the
MAC transcript. The exchange is mandatory: if the ephemeral component is
absent, the session fails closed rather than deriving a non-forward-secret
key.

%=======================================================================
\section{Security Analysis}
\label{sec:proofs}
%=======================================================================
Throughout, $q$ bounds the number of sessions/oracle queries, and reductions
are tight up to the stated additive terms.

\begin{theorem}[Phase-2 mutual authentication]\label{thm:p2}
Under \cref{ass:puf}, EUF-CMA security of $\MAC$, and PRF security of $\KDF$,
\cref{alg:p2} achieves mutual entity authentication with injective agreement
on the transcript $\transcript$. For every PPT $\adv$,
\begin{equation}
\Adv^{\mathrm{p2\text{-}auth}}(\adv)\ \le\
2\,\Adv^{\mathrm{euf}}_{\MAC}(\adv')
+\Adv^{\mathrm{prf}}_{\KDF}(\adv'')
+\frac{q^2}{2^{128}}
+\negl(\secpar).
\end{equation}
\end{theorem}
\begin{proof}
Let $K^{i}_{\mathrm{auth}}=\KDF(\mk_i;\lbl{p2-auth})$. Game $G_0$ is the real
authentication experiment. Game $G_1$ replaces $K^{i}_{\mathrm{auth}}$ by a
uniform key; by PRF security of $\KDF$ and the near-uniformity of $\mk_i$
(\cref{lem:fe}, so $\mk_i$ is unknown to any $\adv$ lacking both the PUF and
the database, by \cref{ass:puf}),
$|\Pr[G_1]-\Pr[G_0]|\le\Adv^{\mathrm{prf}}_{\KDF}+\negl$.

\emph{(GS authenticates UAV.)} A UAV is accepted only if $\sigma_1$ verifies
under $K^{i}_{\mathrm{auth}}$; producing a fresh valid $\sigma_1$ without the
key is an EUF-CMA forgery.

\emph{(UAV authenticates GS.)} Symmetrically for $\sigma_2$: although every
input to $\sigma_2$ is public, the \emph{key} is secret, so a fresh valid tag
is again an EUF-CMA forgery.

A nonce collision that would allow transcript reuse occurs with probability at
most $q^2/2^{128}$. Absent forgery and collision, both accepting parties share
the identical $\transcript$, i.e.\ injective agreement holds. Summing the two
forgery terms gives the bound.
\end{proof}

\begin{corollary}[No challenge read-out oracle]\label{cor:nooracle}
An active $\adv_{\mathrm{chal}}$ impersonating $\GS$ obtains no PUF response
value; hence modelling attacks that require harvested challenge--response
pairs are infeasible, under no assumption beyond \cref{ass:puf}.
\end{corollary}
\begin{proof}
By step~4 of \cref{alg:p2}, the UAV emits $M_3$ only after $\sigma_2$
verifies, which by \cref{thm:p2} an adversary can force with probability at
most $\Adv^{\mathrm{euf}}_{\MAC}+\negl$. When verification fails the UAV emits
nothing; when it succeeds, the wire carries only $\sigma_3$, a MAC over public
transcript fields---never a PUF response or a masked response. No
challenge--response value is therefore exposed.
\end{proof}

\begin{theorem}[Phase-3 peer mutual authentication]\label{thm:p3}
Under EUF-CMA security of $\MAC$ and INT-CTXT security of $\aead$,
\cref{alg:p3} achieves peer mutual authentication with injective agreement,
without GS participation at handshake time and without exposing any per-pair
secret in transit:
\begin{equation}
\Adv^{\mathrm{p3\text{-}auth}}(\adv)\ \le\
2\,\Adv^{\mathrm{euf}}_{\MAC}(\adv')
+2\,\Adv^{\mathrm{int}}_{\aead}(\adv'')
+\frac{q^2}{2^{128}}
+\negl.
\end{equation}
\end{theorem}
\begin{proof}
\emph{Issuance.} $\Cred_{ij}$ is delivered as $\aead.\Enc_{\SK_i}$; by
INT-CTXT the adversary cannot make a UAV accept a different credential, and by
IND-CCA it learns nothing about $\Cred_{ij}$ (used in \cref{thm:secrecy}).
Since $\SK_i$ is a fresh, forward-secret key held only by $\UAV_i$ and the
$\GS$ (\cref{thm:p2,thm:pfs}), the adversary's view is independent of
$\Cred_{ij}$ except with probability $\le\Adv^{\mathrm{int}}_{\aead}$ per
endpoint.

\emph{Handshake.} Given a secret $\Cred_{ij}$, the tokens
$\tau_1,\tau_2,\tau_3$ are MACs under $K^{ij}_{\mathrm{auth}}$; each carries a
distinct type tag and binds both identities and both nonces, so any fresh
acceptance forces a matching sender computation---otherwise an EUF-CMA
forgery. Because credentials are pair-specific, a third party holding some
$\Cred_{k\ell}$ has no information about $\Cred_{ij}$ (independent KDF inputs)
and cannot assert identity $i$ to $j$. Nonce collisions contribute at most
$q^2/2^{128}$.
\end{proof}

\begin{theorem}[Replay and MITM resistance]\label{thm:replay}
\cref{alg:p2,alg:p3} resist replay and active man-in-the-middle attacks.
\end{theorem}
\begin{proof}
Every accepted message is a MAC over a transcript that includes a fresh
$128$-bit nonce, a timestamp validated within $\Delta T$, a per-message type
tag, and---from $M_2$/$P_2$ onward---the peer's nonce. A replay reuses a stale
nonce or timestamp and is rejected by the freshness cache or window; a
cross-session or cross-type substitution alters a MAC input and, by
\cref{thm:p2,thm:p3}, cannot be re-authenticated without the key. Any active
modification invalidates the covering MAC, and injective agreement on
$\transcript$ then guarantees both endpoints hold the identical view. The
distinct type tags further prevent replaying a $P_1$ token as a $P_3$ token.
\end{proof}

\begin{theorem}[Physical-capture resistance]\label{thm:capture}
Capture of $\UAV_k$ (reading its NVM) does not enable $\adv_{\mathrm{cap}}$ to
\textup{(i)}~impersonate any $\UAV_m$ ($m\ne k$) to the $\GS$ in a fresh
Phase~2, nor \textup{(ii)}~impersonate any peer pair $(m,\ell)$ with
$k\notin\{m,\ell\}$ in Phase~3. A captured node's peer-authentication exposure
is exactly its $N-1$ incident links.
\end{theorem}
\begin{proof}
The captured NVM holds $\{\TID_k,P_k,\{\Cred_{kj}\}_j,\ \text{live }\SK\}$.
\emph{(i)}~A fresh Phase~2 as $\UAV_m$ requires
$K^{m}_{\mathrm{auth}}=\KDF(\mk_m;\cdot)$; $\mk_m$ is never stored
(\cref{alg:p1}) and, by \cref{ass:tamper}, is not recoverable from device
$k$'s silicon, while \cref{ass:puf} makes it unpredictable, so forgery reduces
to EUF-CMA. \emph{(ii)}~Phase~3 for $(m,\ell)$ requires
$\Cred_{m\ell}=\KDF(\mk_{\GS};\cdot)$; the captured set $\{\Cred_{kj}\}$ is
independent of $\Cred_{m\ell}$ whenever $k\notin\{m,\ell\}$ (distinct KDF
inputs; $\mk_{\GS}$ never leaves the $\GS$), so acceptance again reduces to
EUF-CMA. Only pairs incident to $k$ share a captured credential, giving the
stated $N-1$ exposure.
\end{proof}

\begin{theorem}[Session-key secrecy]\label{thm:secrecy}
For an adversary that corrupts no endpoint of the target session, $\SK_i$
(resp.\ $\SK_{ij}$) is indistinguishable from uniform:
\begin{equation}
\Adv^{\mathrm{key\text{-}ind}}(\adv)\ \le\
\Adv^{\mathrm{prf}}_{\KDF}+\Adv^{\mathrm{gap\text{-}cdh}}+\negl.
\end{equation}
\end{theorem}
\begin{proof}
By \eqref{eq:sk} the key is $\KDF(k_{\mathrm{lt}}\catt g^{xy}\catt\transcript)$.
Two independent secrets protect it: the long-term key $k_{\mathrm{lt}}$
($\mk_i$ by \cref{lem:fe} and \cref{ass:puf}, or $\Cred_{ij}$ by IND-CCA of
the AEAD delivery, \cref{thm:p3}) and the Diffie--Hellman secret $g^{xy}$.
Modelling $\KDF$ as a PRF, replacing it by a random function costs
$\Adv^{\mathrm{prf}}_{\KDF}$; the input $g^{xy}$ is unpredictable under gap-CDH
given only the authenticated shares $g^{x},g^{y}$ (\cref{thm:p2,thm:p3}). The
output is therefore uniform up to the stated terms.
\end{proof}

\begin{theorem}[Perfect forward secrecy]\label{thm:pfs}
If the ephemeral scalars are erased after use, then revealing the long-term
secrets ($\mk_i$, $\Cred_{ij}$, or even $\mk_{\GS}$) \emph{after} a target
session closes does not compromise that session's key:
\begin{equation}
\Adv^{\mathrm{pfs}}(\adv)\ \le\
\Adv^{\mathrm{gap\text{-}cdh}}+\Adv^{\mathrm{prf}}_{\KDF}+\negl.
\end{equation}
\end{theorem}
\begin{proof}
A post-session \textsf{Reveal} yields $k_{\mathrm{lt}}$ but not the erased
$x,y$. The key~\eqref{eq:sk} still depends on $g^{xy}$, which the adversary
must compute from the transcript shares $g^{x},g^{y}$---a gap-CDH instance,
since those shares were authenticated (no substitution, by
\cref{thm:p2,thm:p3}) and the scalars are gone. Reducing $\KDF$ to a PRF on
the now-unknown input yields indistinguishability up to the stated terms.
Because the ephemeral exchange is mandatory and fail-closed (\cref{sec:p4}),
no non-forward-secret key exists to attack.
\end{proof}

%=======================================================================
\section{Symbolic Verification}
\label{sec:symbolic}
%=======================================================================
We corroborate the authentication and secrecy results with a symbolic model
in the Tamarin prover~\cite{tamarin2013}, whose multiset-rewriting calculus
natively expresses freshness, stateful credential issuance, and compromise
(``reveal'') events.

\paragraph{Model.}
The symbols $h,\MAC,\aead,\KDF$ are functions with the usual equational
theory
\begin{equation}
\aead.\Dec_k\big(\eta;\ \mathit{ad};\ \aead.\Enc_k(\eta;\mathit{ad};m)\big)=m.
\end{equation}
Each PUF is a private per-device function $\mathsf{puf}(\sim\!\mathrm{seed}_i,c)$;
helper data and all wire terms are public. The GS database and issued
credentials are linear facts. Fresh values ($\sim\!N,\sim\!a$) model nonces and
ephemeral scalars; $g^{ab}$ uses the built-in Diffie--Hellman theory.

\paragraph{Trust and compromise rules.}
The Dolev--Yao adversary controls the network. Two additional rules model the
strong adversary: \texttt{Reveal\_NVM(k)} emits $\TID_k,P_k,\{\Cred_{kj}\}$ but
neither $\mathsf{puf}_k$ nor $\mk_k$; \texttt{Reveal\_LT} emits $\mk_i$ or
$\Cred_{ij}$ with an action timestamp used to restrict forward-secrecy queries
to post-session reveals.

\paragraph{Verification lemmas.}
The following lemmas are discharged by \texttt{tamarin-prover -{}-prove} and
correspond one-to-one with \cref{thm:p2,thm:p3,thm:replay,thm:capture,thm:secrecy,thm:pfs}.
\begin{spacing}{1.05}
\begin{verbatim}
lemma p2_injective_agreement:
  "All U GS t #i. Commit(U,GS,<'p2',t>)@i
      ==> (Ex #j. Running(GS,U,<'p2',t>)@j & j<i)
          & not (Ex U2 GS2 #i2. Commit(U2,GS2,<'p2',t>)@i2 & not #i2=#i)"

lemma p3_injective_agreement:
  "All i j t #x. Commit(i,j,<'p3',t>)@x
      ==> (Ex #y. Running(j,i,<'p3',t>)@y & y<x)"

lemma no_replay:
  "All t #a #b. Accept(t)@a & Accept(t)@b ==> #a=#b"

lemma capture_isolation:
  "All m l #x. PeerKey(m,l)@x
      & (Ex k #r. RevealNVM(k)@r & not k=m & not k=l)
      ==> not (Ex #e. K(SKpair(m,l))@e)"

lemma key_secrecy:
  "All k #x. SessionKey(k)@x & not (Ex #r. RevealEndpoint()@r)
      ==> not (Ex #e. K(k)@e)"

lemma forward_secrecy:
  "All k #x. SessionKey(k)@x
      & (All #r. RevealLT()@r ==> x<r)
      ==> not (Ex #e. K(k)@e)"
\end{verbatim}
\end{spacing}

\paragraph{Modelling scope.}
The symbolic layer treats the PUF as a noiseless private function---error
tolerance is a computational correctness property (\cref{lem:fecorrect})
outside the symbolic theory---and abstracts confidential delivery with a
symmetric-encryption operator, leaving the fine-grained confidentiality
argument to the computational analysis (\cref{thm:secrecy}). As a
self-consistency check, a rule that delivered a credential in the clear
falsifies \texttt{key\_secrecy}, confirming that the AEAD wrapping is
load-bearing in the model.

%=======================================================================
\section{Conclusion}
\label{sec:conclusion}
%=======================================================================
We presented a four-phase PUF-based authentication and key-establishment
protocol for UAV swarms. The design roots all authentication in a stable,
PUF-derived master key that is regenerated on demand and never stored, keeps
the raw PUF response entirely on-device, protects peer authentication with
per-pair credentials delivered under authenticated encryption, and makes an
authenticated ephemeral Diffie--Hellman exchange a mandatory component of every
session key. A helper-data privacy lemma with an explicit min-entropy budget
governs parameter selection. Six theorems, each proved by reduction to
standard assumptions and mirrored by a mechanised Tamarin model, establish
mutual authentication in both directions, peer authentication, replay and MITM
resistance, physical-capture resistance, session-key secrecy, and perfect
forward secrecy. Built from a $160$-bit hash, a keyed MAC, a lightweight AEAD,
a KDF, and a single elliptic curve, the protocol is well matched to the
resource constraints of aerial platforms.

%=======================================================================
\begin{thebibliography}{99}
\small
\bibitem{suh2007} G.~E.~Suh and S.~Devadas, ``Physical unclonable functions
for device authentication and secret key generation,'' in \emph{Proc.\ 44th
ACM/IEEE Design Automation Conf.\ (DAC)}, 2007, pp.~9--14.

\bibitem{dodis2008} Y.~Dodis, R.~Ostrovsky, L.~Reyzin, and A.~Smith, ``Fuzzy
extractors: How to generate strong keys from biometrics and other noisy
data,'' \emph{SIAM J.\ Comput.}, vol.~38, no.~1, pp.~97--139, 2008.

\bibitem{krawczyk2010} H.~Krawczyk, ``Cryptographic extraction and key
derivation: The HKDF scheme,'' in \emph{Advances in Cryptology --- CRYPTO},
2010, pp.~631--648.

\bibitem{bernstein2006} D.~J.~Bernstein, ``Curve25519: New Diffie--Hellman
speed records,'' in \emph{Public Key Cryptography --- PKC}, 2006,
pp.~207--228.

\bibitem{bogdanov2011} A.~Bogdanov, M.~Kne\v{z}evi\'c, G.~Leander,
D.~Toz, K.~Var\i c\i, and I.~Verbauwhede, ``SPONGENT: A lightweight hash
function,'' in \emph{Cryptographic Hardware and Embedded Systems --- CHES},
2011, pp.~312--325.

\bibitem{fips202} National Institute of Standards and Technology,
``SHA-3 Standard: Permutation-Based Hash and Extendable-Output Functions,''
FIPS PUB 202, 2015.

\bibitem{ascon} C.~Dobraunig, M.~Eichlseder, F.~Mendel, and M.~Schl\"affer,
``Ascon v1.2: Lightweight authenticated encryption and hashing,''
\emph{J.\ Cryptology}, vol.~34, no.~3, 2021.

\bibitem{bellare1993} M.~Bellare and P.~Rogaway, ``Entity authentication and
key distribution,'' in \emph{Advances in Cryptology --- CRYPTO}, 1993,
pp.~232--249.

\bibitem{dolevyao1983} D.~Dolev and A.~C.~Yao, ``On the security of public key
protocols,'' \emph{IEEE Trans.\ Inf.\ Theory}, vol.~29, no.~2, pp.~198--208,
1983.

\bibitem{tamarin2013} S.~Meier, B.~Schmidt, C.~Cremers, and D.~Basin, ``The
TAMARIN prover for the symbolic analysis of security protocols,'' in
\emph{Computer Aided Verification --- CAV}, 2013, pp.~696--701.
\end{thebibliography}

\end{document}
